\documentclass[11pt,a4paper]{article}

\usepackage{fontspec}
\usepackage{polyglossia}
\setdefaultlanguage{english}
\setmainfont{Liberation Serif}
\setsansfont{DejaVu Sans}
\setmonofont{DejaVu Sans Mono}[Scale=0.82]

\usepackage[margin=2.5cm]{geometry}
\usepackage{amsmath,amssymb}
\usepackage{booktabs}
\usepackage{longtable}
\usepackage{tabularx}
\usepackage{graphicx}
\usepackage{xcolor}
\usepackage{enumitem}
\usepackage{caption}
\captionsetup{font=small,labelfont=bf}
\usepackage{microtype}
\setlength{\emergencystretch}{3em}
\usepackage{pgfplots}
\pgfplotsset{compat=1.18}
\usetikzlibrary{positioning,arrows.meta}
\usepackage[numbers,sort&compress]{natbib}
\usepackage[hidelinks]{hyperref}
\usepackage{url}

\definecolor{skeinc}{HTML}{2F6F9F}
\definecolor{openc}{HTML}{C66A2B}
\definecolor{freshc}{HTML}{2F6F9F}
\definecolor{cachec}{HTML}{9CB8CC}
\definecolor{outc}{HTML}{C66A2B}

\title{\textbf{Skein: A Coding Agent with a Non-Monotone Context}\\[0.4em]
\large Architecture and Measurements on Ported Terminal-Bench Tasks}
\author{Vladimir Lapshin\\[0.3em]
\small Independent researcher\\[0.2em]
\small \texttt{mefrill@yandex.ru}}
\date{Preprint, November 2026}

\begin{document}

\maketitle

\begin{abstract}
A language-model agent reasons probabilistically: its conclusions are plausible but cannot be
proved. Reliability therefore comes not from inside the model but from around it---from a
procedure-discipline. It defines what the model must record, what it may rely on, and when a
piece of work counts as closed. We describe Skein, a coding agent in which this procedure is
built into the journal of work.

The journal stores every action and its result and is never rewritten. The messages for the model
are rebuilt from the journal on every turn, so anything no longer needed can leave the view
without being lost. Two explicit ways of working sit on top of the journal. States fix the goal,
the stopping condition, and the allowed transitions in advance. Dynamic hypotheses let the model
propose checkable claims, work under them, and close them by recording what it actually saw. In
both cases the engine records what was done and on what basis; it does not rely on the model's
own account.

We built the agent and tested it on ported Terminal-Bench tasks. States and hypotheses solve the
same tasks; hypotheses use almost half as many fresh input tokens but take more turns. The main
limitation is the same as for any such agent: the model's decision does not replace verification.
An external check accepts the result, and it is the check that shows where the model went wrong.
\end{abstract}

\noindent\textbf{Keywords:} coding agent, LLM agents, context management, event journal, event
sourcing, context projection, context isolation, prefix caching, work states, Terminal-Bench.

\section{Introduction}

A coding agent is a program that reads a project, makes changes, and checks the result by
following the decisions of a language model. A typical implementation follows the ReAct
loop~\cite{yao2023react}: the model reasons and chooses an action; the executor runs a tool and
returns an observation. The history of such a loop is naturally stored as a tape of messages:
a system instruction, the user's request, and pairs of action and observation. The tape is
monotone---it only grows.

A growing tape has three well-known problems. First, the history exceeds the model's context window on
long tasks. Second, models use long context unevenly: information in the middle is less
accessible than information at the beginning and the end~\cite{liu2023lost}; separately from
this positional effect, we observe an engineering phenomenon---the original goal gets buried deep
into the history (hereafter, goal drift; this is an observation, not a
result~\cite{liu2023lost}). Third, compaction and truncation either lose facts or rewrite them,
and the exact state of the work becomes unrecoverable.

Skein rests on a different assumption: \emph{history is data, and context is a derived
representation}. The engine records actions and observations in an append-only typed event
journal (called IR in the code; hereafter, the journal); the messages for the model are
assembled from the journal on every turn by the function $project()$. This separation is a
direct analogue of event sourcing and command--query responsibility segregation
(CQRS)~\cite{fowler2005event,fowler2011cqrs}: the journal is the only writable model, and the
context is a recomputable read model. The projection is deterministic (a pure function) and
non-monotone in what it shows\footnote{Here ``non-monotone'' means that the visible representation can
shrink when the focus changes; it is not non-monotonic logic, in which conclusions may be
retracted.}: content leaves the context when the focus moves, without disappearing from the
journal.

The second decision is to make the steps of the work explicit and store them in the journal rather
than in code. There are two ways to do this. In states, the steps (phases) are fixed in advance:
each has a goal and a stopping condition, and the allowed transitions are described by the
protocol; the model decides when to move on, and the executor records the transition in the
journal. In dynamic hypotheses, the model chooses the steps as it goes: it proposes a checkable
claim, works under it, and closes it by recording what it saw. In both cases the goal, the
checking condition, and the facts of the work are stored in the journal, so they can be reopened
and checked after the work is done. This is what distinguishes Skein from one-shot subagents and
from graphs hard-coded in the implementation (Section~\ref{sec:states}).

The contributions of this work are:
\begin{enumerate}[leftmargin=1.4em,itemsep=0.2em]
  \item an architecture in which the agent's context is a deterministic non-monotone projection
  of an append-only journal, with addressable access to hidden history (\texttt{recall},
  \texttt{search}); positioning against event sourcing, OpenHands, MemGPT, and subagents;
  \item a state protocol: the goal, stopping criterion, and transition topology are journal
  data interpreted by the protocol rather than by code;
  \item a controlled comparison with opencode on 13 ported Terminal-Bench tasks ($K=5$, two
  network modes), with a decomposition of tokens and cost by state and an analysis of the cost
  structure and the role of prefix caching in the economics of a long tape;
  \item dynamic hypotheses---the next step beyond explicitly defined states---and their controlled
  comparison with states on seven tasks ($K=5$, two network modes), including a long task and a
  task written from scratch;
  \item an analysis of failures and a program for measuring the claimed but as yet unmeasured
  properties (goal retention and controllability).
\end{enumerate}

The source code, documentation, and installation instructions are available in the
repository \url{https://github.com/merfill/skein}, on the website
\url{https://skein-swe.dev}, and as the npm package \texttt{skein-swe}
(\url{https://www.npmjs.com/package/skein-swe}).

The work is built around one question: how to organize a model's reasoning so that it can be
checked after it has been done. The core of Skein is small: an append-only journal, a projection
into messages, and two ways of working on top of the journal---states and dynamic hypotheses. That
is enough to measure the cost of separating context and to show what work with explicit steps
looks like, as opposed to a free-growing tape.

A broader assumption is attached to hypotheses. Models are currently trained mostly on direct
reasoning---the chain of ``think and answer at once''---and such a chain is often unpredictable.
Hypotheses set a different frame: first state a claim, then check it, and only then move on. In
this sense they are a meta-organizer of reasoning. We do not claim that Skein already provides
this organization in finished form; we show how it is built, and we hope that approaches like it
will encourage training models to reason through hypotheses rather than only through direct
chains.

Non-goals: a claim of state-of-the-art performance on a benchmark, cost optimization, and a test
of portability across models. With one API provider and one model, the comparison shows that
Skein is not cheaper than opencode; the goal of the paper is to describe the mechanics and to
measure their consequences.

\section{Related Work}

\subsection{Coding Agents and Phased Pipelines}

SWE-bench~\cite{jimenez2024swebench} made repository-level tasks a standard for agent evaluation; Terminal-Bench~\cite{terminalbench} moves the idea into the terminal. SWE-agent shows
that the design of the agent--computer interface substantially affects
performance~\cite{yang2024sweagent}. Agentless demonstrates that a rigid phased pipeline of
localization, repair, and validation, without autonomous planning, is competitive with
agents~\cite{xia2024agentless}. The state sequence of Skein is close to such a pipeline in
form, but the model chooses transitions between phases according to the protocol rather than a
schedule in code; the goal and criterion of each phase are journal data.

opencode~\cite{opencode} is an open-source terminal agent and our baseline: one growing tape, a
set of tools (shell, read, search, edit, write), and a loop that runs until the model stops or
the limit is exhausted. It represents the main class of implementations.

\subsection{History as Data: Journal, Projection, and Memory}

Event sourcing~\cite{fowler2005event} and CQRS~\cite{fowler2011cqrs} specify a pattern in which
the application state is fully derived from an event journal and the write and read models are
separated. Skein applies this pattern to the context of a language-model agent: the journal is
the source of truth, and the context is a derived representation.

Lost in the Middle~\cite{liu2023lost} motivates abandoning a long tape, but it measures
positional accessibility rather than goal drift. MemGPT~\cite{packer2023memgpt} manages memory
hierarchically, and its main context is also non-monotone. The key difference in Skein is the
way the view changes: MemGPT \emph{edits} the representation through model-driven memory
operations, whereas Skein \emph{recomputes} it from an immutable journal. Editing accumulates
the model's decisions about what to remember; recomputation does not. CoALA~\cite{sumers2024coala}
provides a taxonomy of memory and action spaces; Generative Agents~\cite{park2023generative}
uses a memory stream with retrieval, that is, an editable memory rather than recomputation from
an immutable journal; Letta~\cite{letta} is the production successor of the MemGPT line.
OpenHands~\cite{openhands2024,openhandsSdk2025} is the closest relative: an event stream and a
Condenser that forms the model's representation by discarding events and replacing them with
summaries. OpenHands condensation edits and compresses the representation (the MemGPT line);
Skein recomputes it deterministically from an immutable journal and preserves addressable access
to the hidden part. Our contribution is not event sourcing as such, but its application to the
agent's context together with a state contract in the journal. Reflexion adds verbal feedback on
top of the tape~\cite{shinn2023reflexion}. LangGraph persists graph state through
checkpoints~\cite{langgraphPersistence}, but both the topology and the node contracts are
defined in code.

\subsection{Subagents and Context Isolation}

The orchestrator--workers pattern is described in Anthropic's overview~\cite{anthropic2024building}
and in their multi-agent research system~\cite{anthropic2025multiagent}. In deepagents, a
subagent is an explicit entity: the \texttt{task} tool, an isolated context, its own system
instruction, and a set of tools~\cite{langchain2025deepagents}; the motivation is to avoid
cluttering the parent's context with intermediate results~\cite{dbreunig2025context}.

Functionally, the Skein operations \texttt{general}, \texttt{revision}, \texttt{plan}, and
\texttt{write\_document} are subagents in this sense. The differences in mechanism are: (i)
isolation is structural---each unit of work has its own journal and its own projection rather
than a separate tape; (ii) for states inside \texttt{task}, the goal and stopping criterion are
journal data, and the transition topology is specified by the protocol. The deepagents framework also has
\texttt{CompiledSubAgent}, an arbitrary LangGraph workflow that is technically an automaton as
well; but its goal and criterion are defined by the graph's code rather than by data, and are
therefore not transferable across topologies without changing the code. The subagent
specification in deepagents is declarative too, so the novelty is not declarativeness as such
but the fact that the contract and the transition facts live in the journal and can be checked
after the fact. Declarative frameworks such as CrewAI~\cite{crewai} define roles and tasks in
YAML, but they do not journal transitions and do not recompute the context from an immutable
journal. We return to this distinction in Section~\ref{sec:states}.

\section{Architecture}
\label{sec:arch}

\subsection{Terms}

The terms used in the rest of the paper are listed below.

\begin{table}[h]
\centering
\small
\caption{Terms and their relations.}
\label{tab:terms}
\begin{tabularx}{\textwidth}{@{}lX@{}}
\toprule
Term & Definition \\
\midrule
Session & an ordered set of units of work; the session agent routes messages \\
Unit of work & an isolated journal plus a projection; opened by an operation \\
Operation & the type of a unit of work (\texttt{respond}, \texttt{general}, \ldots; Table~\ref{tab:ops}) \\
State & a protocol node inside a unit of work: goal, criterion, tools, transitions \\
Journal & an append-only sequence of typed events and nodes (IR in code) \\
Projection & the pure function $project(journal) \rightarrow \texttt{Context}$ \\
Turn & one model call (a message $\rightarrow$ action $\rightarrow$ observation cycle);
for opencode we count an agent step \\
\bottomrule
\end{tabularx}
\end{table}

\subsection{Journal and Projection}

The journal is monotone: events are only appended, nodes are never rewritten, and the full
history remains addressable even when it is not shown to the model. A node is a record with the
fields \texttt{id}, \texttt{space}, \texttt{kind}, \texttt{label}, \texttt{payload}, and
\texttt{seq}; an edge has \texttt{id}, \texttt{from}, \texttt{to}, \texttt{kind}, and
\texttt{provenance}. The event types are \texttt{add\_node}, \texttt{add\_edge},
\texttt{descend}, \texttt{return}, \texttt{mutate}, \texttt{revert}, \texttt{rollback}, and
\texttt{record\_rejection}. The node types are \texttt{request}, \texttt{move}, \texttt{turns},
\texttt{action}, \texttt{observation}, \texttt{stop}, \texttt{unactionable},
\texttt{constraint}, \texttt{fact}, \texttt{hint}, \texttt{memo}, \texttt{plan},
\texttt{plan\_item}, \texttt{file}, \texttt{symbol}, and \texttt{test}. The edge types are
\texttt{move}, \texttt{unactionable}, \texttt{turns}, \texttt{turn}, \texttt{stop},
\texttt{result}, \texttt{mutates}, \texttt{plan}, and \texttt{plan\_item}. The events
\texttt{mutate}, \texttt{revert}, and \texttt{rollback} do not rewrite nodes; they record a new
version of an artifact or a rollback, and the versions are kept in the journal. The model's
note in the working instruction is the \texttt{what} field of the \texttt{move} event: it is
written when the state is entered and is not rewritten; the projection reads it as an unverified
guess.

The projection is deterministic: the context is fully determined by the journal and the
projection parameters, which makes it possible to reproduce the context from the journal. Note
that determinism of the projection \emph{does not} mean determinism of the agent's behavior:
the decisions are made by a stochastic model. The property has been checked on 400 random
journals (\texttt{tests/ops/ir\_properties.test.ts}) and on the 125 journals of the measurement campaign (the full series of runs) stored in
the \texttt{result.json} of each cell: two calls to \texttt{project()} on the same journal
produce an identical context (125/125).

The projection is non-monotone. An open state shows:
\begin{itemize}[leftmargin=1.4em,itemsep=0.15em]
  \item \texttt{Start state} and a brief description;
  \item its own turns---a monotone list of actions and observations;
  \item the working instruction for the turn: \texttt{Goal} (the state's goal), the model's
  unverified note, \texttt{Done when} (the stopping criterion), \texttt{Plan} (the state's
  instruction), and an explicit note that the original request is context, not the state's goal.
\end{itemize}
A closed state shows only \texttt{Start state} and \texttt{Finish state}. Closing a state does
not delete its nodes: they remain accessible through \texttt{recall} and \texttt{search}.
\texttt{recall(id, start?, end?)} reads a stored result by the identifier shown on the tape
(from the beginning when no window is given, and paginates a large body when a window is given);
\texttt{search(id, pattern, before?, after?)} matches a regular expression against a stored
result (stdout and stderr) and returns line windows with context. Both tools are read-only, are
always available, and append the returned text to the open state's tape; a repeated
\texttt{recall} without a window is rejected when the body is already in view. There is no
separate token budget: the size is controlled by the window and pagination. The view can shrink
without losing data.

Message 0 is the base role instruction; the user's request is the first \texttt{user} turn. The
state map is not part of the base instruction: the request instruction names the entry
states, and the transition instruction names the reachable states and the rules, locally for the
current move. An open state never sees the goal of another state. The motivation is not to
create a ``second goal'' out of the transition map; this is a design hypothesis, not a measured
result.

\begin{figure}[t]
\centering
\begin{tikzpicture}[
  every node/.style={font=\small},
  box/.style={draw, rounded corners=2pt, align=center, inner sep=5pt, minimum height=8mm},
  arr/.style={-{Stealth[length=2mm]}},
  lbl/.style={font=\scriptsize},
]
\node[box] (journal) at (0,0) {journal\\(append-only)};
\node[box] (proj) at (4.1,0) {$project()$};
\node[box] (ctx) at (7.6,0) {\texttt{Context}\\(view)};
\node[box] (llm) at (11.2,0) {model};
\node[box] (exec) at (7.6,-2.1) {tool\\executor};
\draw[arr] (journal) -- (proj);
\draw[arr] (proj) -- (ctx);
\draw[arr] (ctx) -- (llm);
\draw[arr] (llm) |- node[lbl,right,pos=0.25]{action} (exec);
\draw[arr] (exec) -| node[lbl,below,pos=0.25]{events (append)} (journal);
\draw[arr,dashed] (llm.north) .. controls +(0,1.3) and +(0,1.3) ..
  node[lbl,above]{\texttt{recall} / \texttt{search}} (journal.north);
\end{tikzpicture}
\caption{Work cycle: the journal is the single source of truth; the context is a derived view
recomputed every turn. The dashed arrow is addressable access to hidden history.}
\label{fig:cycle}
\end{figure}

\subsection{Session and Operations}

The session agent parses each user message and either answers it or opens a unit of work with a
self-contained work order (Table~\ref{tab:ops}). Each operation opens its own journal: work on a
document does not carry file reads from an investigation, and fixing a defect does not drag in a
plan.

\begin{table}[h]
\centering
\small
\caption{Session operations.}
\label{tab:ops}
\begin{tabular}{@{}ll@{}}
\toprule
Operation & Purpose \\
\midrule
\texttt{respond} & answer from the dialogue and earlier work \\
\texttt{general} & read-only investigation, no changes \\
\texttt{revision} & a mutating operation already described in the request \\
\texttt{task} & a coding task with a reformulated request \\
\texttt{plan} & a work plan, read-only \\
\texttt{write\_document} & a large Markdown document, section by section \\
\bottomrule
\end{tabular}
\end{table}

\subsection{States: Contract as Data}
\label{sec:states}

A state is defined by four elements: a goal (\texttt{what}), a stopping criterion (\texttt{Done
when}), a set of tools, and admissible transitions. The goal and criterion are journal data that
appear in the working instruction at the decision point. Transitions are specified by the
protocol: the transition instruction names the reachable states and the rules; the model
decides; the executor records the transition fact in the journal.

More precisely, this is an automaton whose \emph{structure} (nodes, contracts, admissible
transitions) is data and whose \emph{transition function} is the model. At present one graph is
implemented inside \texttt{task}; a run may start in \texttt{Patch}, \texttt{Research}, or \texttt{Construction}, and the
transitions are as follows:
\[
\begin{aligned}
&\text{Research} \rightarrow \{\text{End},\ \text{Patch}\}, \\
&\text{Patch} \rightarrow \{\text{Verification},\ \text{Research}\}, \\
&\text{Construction} \rightarrow \{\text{End},\ \text{Research}\}, \\
&\text{Verification} \rightarrow \{\text{End},\ \text{Patch}\}.
\end{aligned}
\]
\texttt{End} is terminal. The architecture permits other topologies but does not implement them; drawing
general conclusions about topologies from one graph would be premature. Note that the model
chooses the transition, so this is not a classical deterministic finite automaton but an
automaton with a data-defined topology and a model as the transition function.

\begin{table}[h]
\centering
\small
\caption{States of the \texttt{task} unit of work: purpose and admissible transitions.}
\label{tab:states}
\begin{tabular}{@{}lll@{}}
\toprule
State & Purpose & Next \\
\midrule
\texttt{Research} & investigate the workspace (read-only) & \texttt{End}, \texttt{Patch} \\
\texttt{Patch} & locate and fix a defect & \texttt{Verification}, \texttt{Research} \\
\texttt{Construction} & create a new artifact from a specification & \texttt{End}, \texttt{Research} \\
\texttt{Verification} & run checks and tests & \texttt{End}, \texttt{Patch} \\
\texttt{End} & finish and report to the parent & --- \\
\bottomrule
\end{tabular}
\end{table}

The difference from a subagent is not that a subagent has no goal (it does: the text of the
subtask), but that the goal, criterion, and transitions live in the journal, while the
transition facts are journaled and can be checked after the fact. ``Data versus code'' is a
consequence, not the main thesis; the difference from \texttt{CompiledSubAgent} and LangGraph is
the same: there the contract is hard-coded in the graph.

The practical consequence we expect is that when the goal and criterion are present at the decision
point, the model is less likely to take the request as a broader task. This is a hypothesis; we
outline how to test it in Section~\ref{sec:future}. The other side is that, as long as the criterion is not
machine-checked, a premature stop is possible---we observed one
(Section~\ref{sec:discussion}).

\subsection{Dynamic Hypotheses: The Next Step Beyond States}
\label{sec:hypotheses}

In states, the work is planned in advance: the phases, their goals, and the transitions are known
before the work starts, and the model only decides when to move on. The next step is to let the
model decide the steps itself, as the work proceeds. For this, a \texttt{task} unit can work with
dynamic hypotheses: the model proposes a checkable claim, works under it, closes it, and moves to
the next one. The way of working is chosen explicitly---\texttt{states} or
\texttt{hypotheses}---and remembered, so that after a pause the work continues in the same way.

A hypothesis is a statement that can be true or false. It is opened by the
\texttt{open\_hypothesis} tool. The model calls the tool with three fields: \texttt{label} (a short
name for the panel), \texttt{claim} (what it is going to establish), and \texttt{done} (what must
be seen for the claim to be confirmed or refuted). The hypothesis number---\texttt{H1},
\texttt{H1.1}, and so on---is assigned by the engine; the model does not write it. The engine
records the opening in the journal. While a hypothesis is open, actions run only under it. A
hypothesis can be simple, with actions under it, or compound, split into sub-hypotheses. Open
hypotheses always form a single chain. The formula container $H$ is created by the engine and is
not written by the model.

A hypothesis is closed not by a tool call but by a text reply with no calls---JSON with the fields
\texttt{verdict}, \texttt{observed}, and \texttt{summary}. In it the model writes whether the
claim was confirmed, refuted, or dropped, what it actually saw, and what the outcome is. When no
hypothesis is open, text closes the goal---JSON with \texttt{outcome} and the final answer. The
mechanism is the same as for states: there the model calls \texttt{move} to enter a state and
closes it with text; here, instead of \texttt{move}, it calls \texttt{open\_hypothesis} and closes
the hypothesis with the same kind of text. \texttt{follows\_from} references show which earlier
hypotheses a step rests on; this is only a reference, not a checked inference.

The messages for the model are assembled as for states but hide more. While a hypothesis is open,
its work is visible. In a closed one, only the opening and closing lines remain; all its turns
leave the messages and stay in the journal, where \texttt{recall} and \texttt{search} can find
them. So the more hypotheses are closed, the shorter the tape, while the journal keeps everything.

The engine checks neither the claim nor the model's decision. It records the opening and the
closing and checks the structure: the fields are non-empty, one chain is open, and the needed
tools are available. The checking condition is written by the model, not by the engine, so its
decision remains an opinion about what it saw. States have the same limitation
(Section~\ref{sec:limitations}), but in a hypothesis it is visible immediately: both the checking
condition and the observation are recorded, so the closing can be checked later.

\section{Experimental Setup}

\subsection{Tasks and Environment}

We use 13 tasks ported from Terminal-Bench~\cite{terminalbench}. Each task runs in a Docker
container: the task places its files, the agent works with the shell and tools, and the task's
built-in verifier assigns a reward (0 or 1). The verifier is created anew from the task cache
at check time, so the reward does not depend on files left by the agent (a trust limitation is
discussed in Section~\ref{sec:limitations}).

Each run is a (task, agent) pair under a fixed network mode, with $K=5$ repeats. The network is
enabled or disabled by an isolating proxy; with \texttt{off}, the network is removed from the
agent only, while the verifier keeps it because its script installs the environment itself. The
factor tests whether the agent can work without external resources. Ten
\texttt{crack-7z-hash} runs were not started with the network off (the task is marked
unsupported in that mode). Both agents use the same model
\texttt{deepseek/deepseek-v4-flash-latest} with reasoning effort set to low, from one provider, through
one API endpoint. The limit is the number of turns, with no wall-clock limit; opencode is
version 1.19.5. The task set and the verifier scripts are the same for both agents, so the
comparison is fair with respect to the model and the checks; the architecture, prompts, and
tools differ (Section~\ref{sec:limitations}).

\begin{table}[t]
\centering
\small
\caption{Experiment protocol.}
\label{tab:protocol}
\begin{tabular}{@{}ll@{}}
\toprule
Parameter & Value \\
\midrule
Tasks & 13 ported Terminal-Bench tasks \\
Repeats & $K=5$ \\
Network modes & on / off (isolating proxy) \\
Agents & Skein; opencode 1.19.5 \\
Model & \texttt{deepseek/deepseek-v4-flash-latest}, reasoning: low; the same for both \\
Concurrency & 10 (\texttt{fix-ocaml-gc}: 3) \\
Limit & 80 model calls; \texttt{fix-ocaml-gc}: 60; +5 in the grace window at \texttt{End} \\
Timeout & none \\
Runs & 260 planned, 250 executed \\
Campaign date & 2026-11-13 \\
\bottomrule
\end{tabular}
\end{table}

\subsection{Metrics and Comparison Sets}

We record the reward, the number of turns, input tokens (fresh and read from cache), output
tokens, reasoning tokens, and cost. The solved rate is computed over all executed runs. Token
comparisons use solved runs (123 per agent). The solved sets \emph{differ in composition} (the
agents fail on different tasks), so the aggregates contain a selection effect. We report both
versions: the full solved sets (Table~\ref{tab:tokens}) and the intersection of 121 pairs in
which both runs are solved (Table~\ref{tab:intersection}). The conclusion holds on the
intersection.

opencode does not return cost. The prices are obtained by solving a 3$\times$3 linear system derived from the provider's
aggregate bill for Skein: fresh
input is 17.80 RUB per 1M tokens, cache read is 0.356 RUB per 1M, so fresh input costs 50$\times$ more than
cache reads, and output is 71.20 RUB per 1M. Cache writes were not charged in these runs
(\texttt{cacheWrite = 0}); reasoning tokens are included in output, so there is no separate
price for them. The 50$\times$ ratio is a price ratio, not a measured effect; the measured cache
share of input is 85.7\% for Skein and 94.9\% for opencode. For opencode the RUB figure is an
estimate under the same price model, not a bill. Prices are in RUB (Russian rubles), the
currency of the provider's bill; we use a reference rate of 80 RUB per USD; at this rate, 178.22 RUB is about
2.23 USD and 125.65 RUB about 1.57 USD. The rate is approximate and was not the billing rate.

\section{Results}

\subsection{Solved Tasks: Parity on a Saturated Suite}

Of the 250 executed runs, 246 are solved (98.4\%); each agent solves 123 of 125. A total of 260
runs were planned (13 tasks $\times$ 2 agents $\times$ 2 network modes $\times$ 5 repeats); 10
\texttt{crack-7z-hash} cells of the task $\times$ agent $\times$ network $\times$ repeat matrix
were not started with the network off, so 250 were executed and
network-mode comparisons use the 12 common tasks. By group: Skein---63/65 (network on) and 60/60
(off); opencode---65/65 and 58/60. Skein's failures are \texttt{fix-git} (one repeat, a
premature stop) and \texttt{fix-ocaml-gc} (one repeat, turn limit), both with the network on.
opencode's failures are two repeats of \texttt{db-wal-recovery} with the network off. Rates: a
premature stop in 1 of 125 Skein runs (0.8\%), a turn-limit miss in 1 of 125 (0.8\%); opencode
has two failures of 125 (1.6\%), both on one task and mode. Across all runs, including failures,
Skein costs 182.72 RUB and opencode about 138.44 RUB; both solve 123 of 125. At the run level,
there are two discordant pairs in each direction (exact McNemar test, $p=1.0$): no difference is detected; the
Wilson 95\% interval for 123/125 is [94.4\%, 99.6\%]. This is not a proof of equivalence but an
absence of detected differences.

An important caveat: our ported subset is saturated (98.4\%) and does not distinguish the
agents by quality; the comparison is informative only about the cost structure. We do not use
external figures for the full Terminal-Bench as a contrast, because benchmark versions and the
leaderboard change. Testing the claimed benefits requires long and difficult tasks
(Section~\ref{sec:future}).

\begin{figure}[t]
\centering
\begin{tikzpicture}
\begin{axis}[
  width=0.95\textwidth, height=6cm,
  ybar, bar width=14pt,
  ymin=0, ymax=70,
  ylabel={runs},
  symbolic x coords={skein_on,skein_off,opencode_on,opencode_off},
  xtick=data,
  xticklabels={Skein (on),Skein (off),opencode (on),opencode (off)},
  enlarge x limits=0.18,
  legend style={at={(0.5,-0.18)},anchor=north,legend columns=3,draw=none,font=\small},
  tick label style={font=\small},
  label style={font=\small},
]
\addplot[fill=skeinc,draw=none] table[x=group,y=solved,col sep=comma]{../figures/fig_solved.csv};
\addplot[fill=outc,draw=none] table[x=group,y=failed,col sep=comma]{../figures/fig_solved.csv};
\addplot[fill=cachec,draw=none] table[x=group,y=unsupported,col sep=comma]{../figures/fig_solved.csv};
\legend{solved,failed,not run}
\end{axis}
\end{tikzpicture}
\caption{Outcomes by group. There are 260 runs in total; 10 were not run
(\texttt{crack-7z-hash}, network off).}
\label{fig:solved}
\end{figure}

\subsection{What the Journals Show}

Across 125 Skein runs, \texttt{recall} was never called and \texttt{search} was called twice in
two runs; in the same runs \texttt{bash} was called 1985 times, \texttt{read} 170 times, and
\texttt{grep} 8 times. Thus addressable access to hidden history was almost unused on these
short tasks, and the benefits of \texttt{recall}/\texttt{search} are not tested here. States
were used moderately: 301 \texttt{move} events across 125 runs (median 2); the engine rejected 22
actions out of 2518 (0.9\%).

\subsection{Tokens and Cost: Less Volume, a Higher Bill}

opencode uses 1.49$\times$ more total tokens, but almost all of the excess is cache read (94.9\%
of input vs 85.7\% for Skein). Skein is more compact in total input, but it has 1.87$\times$ more
fresh input tokens and 1.33$\times$ more output tokens: units of work do not share a growing
prefix and may re-read the workspace. Skein also has more turns (1.36$\times$). With a
50$\times$ price difference between cache and fresh input, this makes opencode cheaper: 125.65
vs 178.22 RUB, that is, opencode is 29.5\% cheaper and Skein 1.42$\times$ more expensive
(Table~\ref{tab:tokens}, Figure~\ref{fig:tokens}).

The cache share in the first repeat is close to that in later repeats (Skein 86.7\% vs 85.3\%;
opencode 94.4\% vs 95.2\%), so cross-run prefix reuse does not explain the gap. At the task
level, opencode is cheaper on 12 of 13 tasks (sign test, $p\approx0.0034$); it is more expensive
only on \texttt{fix-ocaml-gc}.

Decomposition: Skein makes 1.36$\times$ as many steps, but only 1.27$\times$ as many model API
calls. Output
per turn is nearly equal (507 vs 520), while fresh input per API call is 1.47$\times$ larger
(1.37$\times$ per step). \texttt{End} accounts for 9.8\% of turns; without \texttt{End} the step ratio is 1.22.
The growth in fresh input therefore comes from more calls at almost the same output per call, not
longer answers.

Tool counts over 125 runs: Skein---1985 \texttt{bash}, 170 \texttt{read}, 113 \texttt{write},
70 \texttt{edit}, 48 \texttt{list}, 8 \texttt{grep}; opencode---1982 \texttt{bash}, 75
\texttt{read}, 83 \texttt{write}, 70 \texttt{edit} (it has no \texttt{grep}/\texttt{list};
search goes through \texttt{bash}). Skein makes more explicit \texttt{read}s (170 vs 75), while
\texttt{bash} is nearly the same; this does not prove repeated reading across states
(Section~\ref{sec:discussion}).

By API calls: Skein has 2447 (one turn is one call), opencode has 1923 calls at 1800 steps (1.07
calls per step). On solved runs the ratio is 1.27, not 1.36 by steps. Fresh input per API call
is 1840 vs 1252 tokens (1.47$\times$; 1.37$\times$ per step). The number of \texttt{bash} calls
is nearly equal (1985 vs 1982), so the difference in turns reflects granularity and protocol overhead, not a
larger amount of terminal work. A Skein turn is one model call and may contain several tool actions; the
action counts therefore exceed the turn count, and a \texttt{move} is an action inside a turn, not an
extra turn.

\begin{table}[t]
\centering
\small
\caption{Tokens and cost, solved runs only (123 per agent; the sets differ in composition,
see \S4.2). Ratios use unrounded values. Turns: for Skein, model calls; for opencode, agent
steps; opencode's model API calls are a separate row. opencode cost is an estimate.}
\label{tab:tokens}
\begin{tabular}{@{}lrrr@{}}
\toprule
 & Skein & opencode & opencode/Skein \\
\midrule
Input, total (M) & 31.38 & 47.63 & 1.52 \\
\quad cache read (M) & 26.88 & 45.22 & 1.68 \\
\quad fresh (M) & 4.50 & 2.41 & 0.53 \\
Output (M) & 1.24 & 0.94 & 0.75 \\
Reasoning (M) & 0.81 & 0.62 & 0.77 \\
Total, input+output (M) & 32.62 & 48.56 & 1.49 \\
Cache share of input & 85.7\% & 94.9\% & --- \\
Turns/steps & 2447 & 1800 & 0.74 \\
Model API calls & 2447 & 1923 & 0.79 \\
Cost (RUB) & 178.22 & \textasciitilde125.65 & 0.71 \\
\bottomrule
\end{tabular}
\end{table}

\begin{table}[t]
\centering
\small
\caption{Intersection check: 121 run pairs in which both agents solve the task. opencode cost
is an estimate. Ratios are opencode/Skein.}
\label{tab:intersection}
\begin{tabular}{@{}lrrr@{}}
\toprule
 & Skein & opencode & opencode/Skein \\
\midrule
Input, total (M) & 31.08 & 45.11 & 1.45 \\
\quad cache read (M) & 26.64 & 42.79 & 1.61 \\
\quad fresh (M) & 4.44 & 2.32 & 0.52 \\
Output (M) & 1.22 & 0.92 & 0.76 \\
Reasoning (M) & 0.80 & 0.62 & 0.78 \\
Turns/steps & 2407 & 1751 & 0.73 \\
Cost (RUB) & 175.60 & 122.26 & 0.70 \\
\bottomrule
\end{tabular}
\end{table}

On the intersection, Skein is 1.44$\times$ more expensive (opencode is 30.4\% cheaper), so the
conclusion does not change. A paired bootstrap over the 121 pairs (5000 resamples) gives a 95\%
interval for the opencode/Skein cost ratio of 0.70 [0.59, 0.81]; a cluster bootstrap over the 13
tasks gives [0.50, 0.85]. For the fresh-input ratio the interval is 0.52 [0.44, 0.61]. All three
intervals exclude 1, so the direction is robust to selection and clustering. Pairing by repeat
index is artificial (runs are independent), and the coverage of a cluster interval over 13 tasks
is unstable; the intervals are therefore descriptive rather than a strict paired design.

\begin{figure}[t]
\centering
\begin{tikzpicture}
\begin{axis}[
  width=0.8\textwidth, height=6cm,
  ybar stacked, bar width=26pt,
  ymin=0, ymax=55,
  ylabel={M tokens},
  symbolic x coords={Skein,opencode},
  xtick=data,
  enlarge x limits=0.5,
  legend style={at={(0.5,-0.16)},anchor=north,legend columns=3,draw=none,font=\small},
  tick label style={font=\small},
  label style={font=\small},
]
\addplot[fill=freshc,draw=none] table[x=agent,y=fresh,col sep=comma]{../figures/fig_tokens.csv};
\addplot[fill=cachec,draw=none] table[x=agent,y=cache,col sep=comma]{../figures/fig_tokens.csv};
\addplot[fill=outc,draw=none] table[x=agent,y=output,col sep=comma]{../figures/fig_tokens.csv};
\legend{fresh input,cache,output}
\end{axis}
\end{tikzpicture}
\caption{Token composition, solved runs (aggregates over 123 runs per agent). opencode's extra
volume is almost entirely cache; Skein has a larger fresh-input share.}
\label{fig:tokens}
\end{figure}

The medians and interquartile ranges per run are given in Table~\ref{tab:medians}. The median
opencode run is cheaper; Skein does not aim at minimum cost, so the median cost says little
about the quality of context management.

\begin{table}[t]
\centering
\small
\caption{Medians per run and interquartile ranges (in brackets) on the 12 common tasks (without
\texttt{crack-7z-hash}). Turns: for Skein, model calls; for opencode, steps. opencode cost is an
estimate.}
\label{tab:medians}
\begin{tabular}{@{}lrrrr@{}}
\toprule
Agent / network & Turns/steps & Fresh input, k & Output, k & Cost, RUB \\
\midrule
Skein, on & 14 [11--30] & 22.2 [17.7--56.8] & 4.5 [2.8--14.5] & 0.70 [0.54--2.25] \\
Skein, off & 15 [10.8--24.3] & 25.0 [17.4--46.7] & 7.3 [3.0--12.3] & 1.00 [0.58--1.99] \\
opencode, on & 9 [7--20] & 10.3 [4.1--26.5] & 2.9 [1.6--10.1] & 0.40 [0.20--1.42] \\
opencode, off & 9 [7--19] & 10.4 [4.7--26.4] & 2.7 [1.5--11.4] & 0.44 [0.23--1.37] \\
\bottomrule
\end{tabular}
\end{table}

On the 12 common tasks the anomaly persists: the median output of Skein is higher with
\texttt{off} than with \texttt{on} (7.3k vs 4.5k), while for opencode it is lower (2.7k vs
2.9k); this is not an effect of task composition. The 1.42$\times$ difference is based on sums; the
medians are 0.70 vs 0.40 (a ratio of 1.75), but the IQRs overlap substantially, and Skein is
cheaper in 23 of the 121 pairs (19\%). The direction is stable, but the magnitude depends on
the aggregate.

\subsection{Decomposition by Skein State}

Across all 125 Skein runs, \texttt{Construction} is the most expensive state (35.5\% of cost),
followed by \texttt{Patch} (26.1\%) and \texttt{Research} (23.8\%). \texttt{Verification} is
cheap: 8.6\% of cost at 11.6\% of turns. \texttt{Patch} and \texttt{Verification} together
account for 34.7\% of cost and 36.7\% of turns (925 of 2518). \texttt{End} takes about two turns
per run (248/125) and 6.0\% of cost---the overhead of finishing. \texttt{Construction} is
expensive mainly because of \texttt{cobol-modernization} (32.7 RUB) and
\texttt{regex-log} (14.0 RUB), tasks that create an artifact from scratch; this is task
composition, not inefficiency of the state itself. The decomposition is approximate: the model
chooses the state label, and the transition turn is assigned to the new state.

\begin{table}[t]
\centering
\small
\caption{Cost by Skein state (all 125 runs, including failures).}
\label{tab:state}
\begin{tabular}{@{}lrrr@{}}
\toprule
State & Turns & Turn share & Cost \\
\midrule
\texttt{Construction} & 760 & 30.2\% & 64.86 RUB (35.5\%) \\
\texttt{Patch} & 634 & 25.2\% & 47.67 RUB (26.1\%) \\
\texttt{Research} & 585 & 23.2\% & 43.49 RUB (23.8\%) \\
\texttt{Verification} & 291 & 11.6\% & 15.71 RUB (8.6\%) \\
\texttt{End} & 248 & 9.8\% & 10.98 RUB (6.0\%) \\
\bottomrule
\end{tabular}
\end{table}

\subsection{Hypotheses versus States}
\label{sec:hyp-exp}

We tested hypotheses on seven tasks where a checkable claim is natural: finding a cause and
recovery (\texttt{fix-git}, \texttt{git-leak-recovery}, \texttt{custom-memory-heap-crash},
\texttt{password-recovery}, \texttt{db-wal-recovery}, \texttt{fix-code-vulnerability},
\texttt{fix-ocaml-gc}). The conditions are the same: $K=5$, network on and off, one model. Both
ways of working run in one series, so solved tasks, turns, and fresh input are comparable; RUB
costs are comparable only within this series, because the day's price differs from the states
series. The limit is each task's own limit; on \texttt{fix-ocaml-gc} the hypotheses were also run
with an 80-turn limit. The task written from scratch is \texttt{modernize-scientific-stack}.

\begin{table}[t]
\centering
\small
\caption{States and hypotheses on seven tasks ($K=5$, network on/off). Medians per run; cost is at
the series-day price, comparable only within the table.}
\label{tab:hyp}
\begin{tabular}{@{}lrrrr@{}}
\toprule
Way of working / network & Solved & Turns & Fresh input, k & Cost, RUB \\
\midrule
States, on & 35/35 & 17 & 34.5 & 2.42 \\
States, off & 34/35 & 15 & 28.7 & 1.61 \\
Hypotheses, on & 34/35 & 20 & 19.6 & 1.88 \\
Hypotheses, off & 34/35 & 22 & 20.1 & 1.98 \\
\bottomrule
\end{tabular}
\end{table}

Hypotheses solve as many tasks as states: 34--35 of 35 in each group. The only failure is a lack
of turns on the long task. The main difference is in the messages the model sees: hypotheses use
almost half as many fresh input tokens (19.6--20.1k vs 28.7--34.5k per run), because the work of a
closed hypothesis disappears from the tape. This is paid for in turns: 20--22 vs 15--17; opening
and closing a hypothesis take separate steps. The cost picture is mixed: on short tasks
hypotheses are cheaper (\texttt{fix-git} 0.79/0.98 vs 1.43/1.30 RUB, \texttt{custom-memory}
3.2/3.5 vs 3.9/5.0, \texttt{fix-code} 1.6/0.83 vs 2.9/1.6), but on \texttt{db-wal-recovery}
without the network and on the long task they cost more in turns.

\texttt{fix-ocaml-gc} shows the boundary of the way of working. At a 60-turn limit hypotheses
solve 8 of 10 (both failures are \texttt{max\_turns}), states 9 of 10. At 80 turns hypotheses
solve 10 of 10 (medians 44 and 68 turns, 1.98 and 3.42 RUB). So the failure at 60 comes from the
budget, not from the way of working: a long task needs spare turns.

The control task \texttt{modernize-scientific-stack}: hypotheses solve 9 of 10 (18 turns, about
0.58 RUB). The single failure is instructive. The script computed the mean correctly
($-15.5^\circ$C) but printed an extra line with a range before it; the check looks for the first
number after ``Station 101'' and takes $-4.3$. The computation is correct, but the answer was not
accepted. This is not a property of hypotheses but a general limitation
(Section~\ref{sec:limitations}): the model does not decide whether the result is accepted, and the
check also looks at the format. With hypotheses this is more visible: the model recorded that the
goal was reached, but someone else accepted the result.

\subsection{What the Hypotheses Journals Show}
\label{sec:hyp-journals}

Hypotheses are not just a declaration: they are actually created and closed. In the 70 runs of the
hypotheses way of working, every run has at least one hypothesis; there are 99 in total, most
often one per run (47 runs), two in 19, and three or four in four. Of these, 95 are closed: 93
confirmed, 2 refuted, 0 dropped. Another 4 remained open, all in four \texttt{fix-ocaml-gc} runs
that hit the 60-turn limit; two of those runs passed the check anyway. Nested hypotheses
(\texttt{H1.1}) are rare: only 4; \texttt{follows\_from} references are filled in for 21 of the 99
hypotheses.

Hypotheses reflect the progress of the solution. In each of the 66 runs where the goal is closed,
it is closed after the last closed hypothesis: while one is open, the work is not considered
finished. More hypotheses mean more turns: runs with one hypothesis have a median of 20 turns,
runs with two or more have 28. On the long \texttt{fix-ocaml-gc}, raising the budget from 60 to 80
turns increases both the number of hypotheses (16 $\rightarrow$ 21) and the number of closed goals
(6 $\rightarrow$ 9 of 10): the extra turns go to new checkable steps, not to repeating what was
already done.

A closed goal does not mean that the result is accepted. In the control task \texttt{modernize},
all 10 runs closed the goal as \texttt{achieved}, but the check accepted 9: in one case the
computation was correct but the answer format was not. The journal shows that the model considers
the work finished; the external check still decides whether the result is accepted.

\section{Discussion}
\label{sec:discussion}

\subsection{The Economics of Cache and Output}

Fresh input costs 50$\times$ more than cache reads (17.80 / 0.356), and the
measured cache share of input is 85.7\% for Skein and 94.9\% for opencode. An agent with a larger
total volume (opencode) therefore turns out to be cheaper than an agent with a smaller volume
but a larger fresh-input share (Skein).

The 52.6 RUB gap decomposes not only into cache (Table~\ref{tab:costsplit}): fresh input
contributes +37.2 RUB, output +21.6 RUB, and cache saves 6.5 RUB. Output is the largest cost
item for both agents (49.5\% for Skein, 53.1\% for opencode), and about 66\% of output is
reasoning. Without cache, Skein is 29\% cheaper (646.9 vs 914.5 RUB): cache reverses the sign (at
the current price and cache share). If Skein had opencode's cache share (94.9\%) at its own
input, the costs would be nearly equal (127.4 vs 125.7 RUB). The optimization target is
therefore more precise: stabilizing the prefix inside a state, rather than reducing the number
of tokens as such.

Price sensitivity (Table~\ref{tab:sensitivity}): at a fresh/cache ratio of 5$\times$, Skein is
slightly cheaper (264.1 vs 270.6 RUB); at 10$\times$, opencode is 12\% cheaper; at 50$\times$,
29\%; at 100$\times$, 32\%. Break-even is about 5.5$\times$. To match opencode at the current
price, Skein would need a cache share of 95.2\% (it is now 85.7\%).

\begin{table}[t]
\centering
\small
\caption{Cost at different fresh/cache price ratios (solved runs; token volumes are fixed to the
measured values and only the prices are recomputed). At 5$\times$ Skein is slightly cheaper;
break-even is about 5.5$\times$.}
\label{tab:sensitivity}
\begin{tabular}{@{}rrrr@{}}
\toprule
fresh/cache & Skein, RUB & opencode, RUB & opencode/Skein \\
\midrule
5   & 264.1 & 270.6 & 1.03 \\
10  & 216.2 & 190.1 & 0.88 \\
20  & 192.3 & 149.8 & 0.78 \\
50  & 178.0 & 125.7 & 0.71 \\
100 & 173.2 & 117.6 & 0.68 \\
\bottomrule
\end{tabular}
\end{table}

The practical conclusion: under the current caching regime, total tokens are a poor proxy for
cost, and cost is a poor proxy for the quality of context management. Comparing architectures
by money on short tasks is uninformative in both directions.

\begin{table}[t]
\centering
\small
\caption{Decomposition of the cost gap (Skein $-$ opencode, solved runs). The sum differs from
52.6 RUB because of rounding.}
\label{tab:costsplit}
\begin{tabular}{@{}lrr@{}}
\toprule
Component & Contribution, RUB & Share of gap \\
\midrule
Fresh input ($+2.09$M $\times$ 17.80) & $+37.2$ & $+71\%$ \\
Output ($+0.30$M $\times$ 71.20) & $+21.6$ & $+41\%$ \\
Cache ($-18.34$M $\times$ 0.356) & $-6.5$ & $-12\%$ \\
\midrule
Total & $+52.3$ & --- \\
\bottomrule
\end{tabular}
\end{table}

\subsection{The Cost of the Current Implementation: Prefix and Re-reading}

Context isolation is not free: the brief summary of the previous unit does not carry the content
that was read, so the next unit may read the workspace anew; on a tape, facts are reused from
the cache. The measurement is possible only for explicit \texttt{read}s: there are 170 \texttt{read}
calls, 9 of which (5.3\%) target a file already read in a different state; reads made through
\texttt{bash} (1985 calls) are not separately accounted for. The first working turn of a state has a cache share
of 66.2\% versus 87.1\% on subsequent turns; this is direct evidence of prefix instability when
entering a state, although the first turns account for only about 2\% of input. Entry effects alone explain
only about 0.4 percentage points of the 9.2-point cache-share gap; the remainder is within-state
prefix behavior. We interpret the
increase in fresh input as the cost of the current implementation of isolation---prefix
instability and more turns---not as a proven property of the architecture. An illustrative case
(not part of the experiment: a different setup, one repository, six user messages, $K=1$): Skein
took 63 steps versus 24 for opencode, 113.7k fresh input versus 35.8k, 3.71 RUB versus about
1.46 RUB, and a cache share of 67\% versus 95\% in a short session. Whether there is a gain on
long work is an open question. Prediction: for a monolithic tape, compaction or truncation may
invalidate the prefix, and then opencode's cache advantage disappears precisely on long tasks;
this is a hypothesis, not a measurement. The two longest tasks by total turns are
\texttt{fix-ocaml-gc} (392) and \texttt{custom-memory-heap-crash} (390). On the first, Skein is
cheaper per solved run (2.58 vs 2.81 RUB); on the second it is more expensive (2.87 vs 1.38
RUB); this is more likely a selection effect than a property of the architecture.

\subsection{States and Subagents: What Exactly Differs}

Functionally, \texttt{general}, \texttt{revision}, \texttt{plan}, and \texttt{write\_document}
are subagents. The substantive difference of \texttt{task} is the contract as data: the goal,
criterion, and transition topology live in the journal and are interpreted by the protocol, and
transition facts are recorded. This makes topologies objects of study (they can be changed,
compared, and reproduced without changing code) rather than a property of the implementation.
At present a single state graph is implemented; the question of other topologies (migrations,
incident investigations, a document with fact checking) remains open---a separate direction,
not the current phase of the project.

\subsection{States and Hypotheses: Who Splits the Work into Steps}

Both ways use one journal and the same messages. The difference is who splits the work into steps.
States do this in advance: the phases, their goals, and the transitions are known before the work
starts, and the model only decides when to move on. In hypotheses the model decides this itself as
the work proceeds: it proposes checkable claims, and the messages show only the work of the open
hypothesis. States therefore fit tasks with a ready specification or procedure, while hypotheses
fit tasks where a cause must be found and something must be checked. The task written from scratch
confirms this only partly: hypotheses solve it, but they take more turns and are sensitive to the
result format.

The common part is the checking condition at the moment of decision. For a state it is
\texttt{Done when}; for a hypothesis it is \texttt{done} plus the recorded \texttt{observed}. In
both cases the engine does not check the condition; it records the model's opinion. A model's
decision is therefore not truth: it makes a step checkable later, but does not prove it. The
difference is that a hypothesis has a condition and an observation at every step, while a state
has them only at its boundary.

\subsection{Failures}

\texttt{fix-git}, Skein, network on, r4: the agent found a dangling commit but did not merge it
into the branch; the \texttt{Research} state closed with a textual diagnosis, after which the
model moved to \texttt{End}. The \texttt{Research}$\rightarrow$\texttt{End} transition is allowed
by the topology, but its precondition,
``the required result exists,'' was not satisfied. The engine does not check the criterion, so
the transition was recorded as admissible. This is a premature stop, not an inadmissible
transition. The other repeats performed the merge.

\texttt{fix-ocaml-gc}, Skein, network on, r5: the limit of 60 turns was exhausted; as the run
approached \texttt{End}, the engine granted a grace window of 5 more turns, so the run stopped at 65. The
verifier did not find the expected test output. This is a turn-limit miss, not a harness
failure.

\texttt{db-wal-recovery}, opencode, network off, r2 and r4: 40--57 steps, 2.3M--4.2M input
tokens (mostly cache). The cause has not been analyzed; this is material for analysis, not a
conclusion about the method.

\section{Limitations}
\label{sec:limitations}

\begin{itemize}[leftmargin=1.4em,itemsep=0.2em]
  \item \textbf{Saturated suite.} 98.4\% solved: the subset does not distinguish quality; the
  conclusions are limited to the cost structure on short tasks.
  \item \textbf{One model, one API provider.} Reasoning effort set to low; portability was not tested.
  The exact model snapshot is not pinned: the name contains \texttt{latest}.
  \item \textbf{Different tools and prompts.} The agents' tool sets differ; no
  \texttt{recall}/\texttt{search} ablation was performed, so the effect of the architecture is
  not separated from the effect of the tool interface.
  \item \textbf{opencode cost is an estimate} based on prices derived from Skein's bill. Token
  conclusions are more reliable than monetary ones.
  \item \textbf{Different solved sets.} Aggregates over 123 solved runs contain a selection
  effect; the intersection is 121 pairs, and its results are reported in
  Table~\ref{tab:intersection}.
  \item \textbf{Small sample.} $K=5$, 13 tasks; medians and interquartile ranges are reported,
  and bootstrap intervals for ratios are computed, but there are no full cluster-level tests.
  \item \textbf{Contamination.} The Terminal-Bench tasks are public and may have been in the
  model's training data; we did not control for this.
  \item \textbf{The stopping criterion is not machine-checked.} A state closes with a
  textual reply from the model; cycle and premature-stop detection is deferred.
  \item \textbf{Determinism of the projection $\neq$ reproducibility of behavior.} The projection
  was checked on 400 synthetic and 125 real journals (identical context); the agent's behavior
  is stochastic.
  \item \textbf{The key benefits are not measured.} That the agent keeps its goal and can be
  controlled is still an assumption, not a result.
  \item \textbf{Repeated reads are not measured directly.} In the Skein journals,
  \texttt{read} appears 170 times, 9 of which (5.3\%) re-read a file in another state; most work
  goes through \texttt{bash}, so we do not isolate the cost of re-reading.
  \item \textbf{One state topology}; the state accounting is approximate (the model chooses the
  label).
  \item \textbf{Trust in the verifier.} The verifier scripts are staged after the agent, but the
  agent could in principle have changed system binaries; verifier isolation was not hardened.
  \item \textbf{Prompt injection.} \texttt{recall}/\texttt{search} return untrusted file text to
  the tape; for an agent working on third-party repositories this is an attack surface, and we
  did not defend against it. For third-party documents (contracts, for example) it is also a risk
  of taking text for an instruction.
  \item \textbf{Sampling parameters are not pinned.} Temperature and related parameters are not
  specified in the protocol.
  \item \textbf{Task selection.} The 13 tasks are the available ported Terminal-Bench tasks, not
  a random sample; external validity is limited.
  \item \textbf{Projection cost.} \texttt{project()} recomputes the context over the whole
  journal on every turn ($O(n)$ per turn); long tasks will need snapshots or memoization.
  \item \textbf{The interaction with the network mode is unexplained.} The median output of
  Skein is higher with \texttt{off} than with \texttt{on} (7.3k vs 4.5k), while for opencode it
  is lower; we do not interpret this.
  \item \textbf{Environment.} Only \texttt{/bin/bash} (Windows: WSL2/Docker); 13 ported tasks,
  not the full set.
  \item \textbf{The hypothesis decision is not checked.} The engine checks only the closing
  structure (non-empty fields, one open chain); the claim and the decision are the model's opinion.
  \item \textbf{The hypotheses series is smaller.} Seven tasks, $K=5$, one model; the task written
  from scratch is one task. RUB costs for hypotheses are comparable only within this series.
  \item \textbf{The way of working is chosen explicitly.} The system does not decide on its own
  where hypotheses are needed and where states are; the user chooses.
  \item \textbf{\texttt{follows\_from} is not checked.} References are stored, but the engine does
  not check that the conclusion really follows from the cited hypotheses.
  \item \textbf{A long task needs more turns.} On \texttt{fix-ocaml-gc} hypotheses need an 80-turn
  budget; at 60 they hit the limit although the task is solvable.
  \item \textbf{The check also looks at the format.} On \texttt{modernize} a correct computation
  failed because of an extra line in the output: the check is format-sensitive, and the model's
  decision does not change that.
\end{itemize}

\section{Availability}

The source code, documentation, and campaign materials are available at
\url{https://github.com/merfill/skein}. A project description and installation instructions are
at \url{https://skein-swe.dev}. The npm package is \texttt{skein-swe}
(\url{https://www.npmjs.com/package/skein-swe}); the Docker image is
\url{https://hub.docker.com/r/mefrill/skein-swe} (the Docker Hub account is \texttt{mefrill},
not \texttt{merfill}). Documentation is in \texttt{docs/en} and \texttt{docs/ru}; campaign
materials are in \texttt{bench/results/campaign-k5/}; the cell journals are in
\texttt{result.json} (field \texttt{events}), with a 125/125 replay; raw runs are published with
the release. The text is licensed under CC BY 4.0.

\section{Conclusion and Future Work}
\label{sec:future}

What is measured: on a saturated suite of short tasks, Skein and opencode solve the same number;
Skein uses fewer total tokens but more fresh input and output, and at a 50$\times$ price
difference for cache it is 1.42$\times$ more expensive. We interpret this as the cost of the
current implementation of isolation: prefix instability and more turns; with opencode's cache
share, the costs are nearly equal (127.4 vs 125.7 RUB). Goal drift was not measured; explicit
repeated \texttt{read}s are rare (9 of 170), and reads through \texttt{bash} are not separately accounted for.
We have not yet measured whether the agent keeps its goal at the decision point or how
controllable it is.

The next phase is the development of the journal as the basis of context, not an expansion of
functionality for its own sake:
\begin{enumerate}[leftmargin=1.4em,itemsep=0.2em]
  \item non-monotone context: the first step is already done---dynamic hypotheses, where the work
  of a closed hypothesis leaves the tape and stays in the journal; next are automatic choice of
  the way of working and machine-checked conditions: the model states \texttt{done} as a formal
  specification (a predicate, a regular expression, an AST pattern), and the engine checks it
  before closing; the same for \texttt{follows\_from} references;
  \item reasoning (the Ankyra project---a research project in which the journal is used as a
  store for reasoning traces): the journal should store not only actions and observations but
  also reasoning traces, and the projection should choose which of them to show;
  \item rollback to checkpoints: the journal allows returning to a saved state of the work and
  continuing from another branch; this is an operation on the journal, not on the tape;
  \item the observer agent: it is currently non-operational; restoring it is again work on the
  journal, because the observer reads the journal and proposes hints or a stop.
  \item a full hypotheses campaign on 13 tasks and a second model: the current series covers seven
  tasks and one model, and the domain of building a program from scratch is measured on one task.
\end{enumerate}

An architecture with an immutable journal and checkable steps is especially apt where an audit
trail is needed: in legal and formal texts, where it matters not only to get an answer but to show
where it came from.

\appendix

\section{Reproduction}

\begin{sloppypar}
Materials: \texttt{bench/results/campaign-k5/}---the protocol (\texttt{manifest.json}), summaries
(\texttt{summary.md}, \texttt{summary.json}, \texttt{summary.csv}), token comparison
(\texttt{token-comparison.md}), and per-state cost (\texttt{state-cost.md}). Raw runs and
journals are in \texttt{bench/runs/matrix/campaign-k5/}; the cell journals are in
\texttt{*/*cells/*/result.json} (field \texttt{events}), with a 125/125 replay.
\end{sloppypar}

The campaign is run through \texttt{tests/sandbox/campaign.ts}: \texttt{--repeat 5},
\texttt{--concurrency 10}, \texttt{--task-concurrency fix-ocaml-gc=3}, both network modes, and no
wall-clock limit; a trial run without containers is \texttt{--dry}. Docker, the task cache, and
the opencode binary are required; \texttt{SKEIN\_LIVE=true}.

The series with hypotheses uses the same \texttt{campaign.ts}/\texttt{matrix.ts} with
\texttt{--agents skein,skein-hypotheses}; the mode is selected by \texttt{--mode hypotheses}
(harness) or \texttt{SKEIN\_TASK\_MODE=hypotheses} (agent). Artifacts:
\texttt{bench/runs/matrix/inquiry-v26-modes}, \texttt{inquiry-v27-ocaml-80},
\texttt{inquiry-v28-modernize}.

\section{Per-Task Results}

\begin{table}[h]
\centering
\small
\caption{Solved runs and cost by task (solved only; opencode cost is an estimate). The
denominators differ, so normalized RUB/run columns are added. For \texttt{crack-7z-hash} the
denominator is 5 (network off was not run). Sums may differ from Table~\ref{tab:tokens} by a few
slightly, because of rounding.}
\label{tab:pertask}
\begin{tabular}{@{}lrrrrrr@{}}
\toprule
Task & Skein & opencode & S, RUB & O, RUB & S, RUB/run & O, RUB/run \\
\midrule
cobol-modernization & 10/10 & 10/10 & 36.73 & 34.09 & 3.67 & 3.41 \\
crack-7z-hash & 5/5 & 5/5 & 5.32 & 4.98 & 1.06 & 1.00 \\
custom-memory-heap-crash & 10/10 & 10/10 & 28.68 & 13.78 & 2.87 & 1.38 \\
db-wal-recovery & 10/10 & 8/10 & 8.64 & 5.95 & 0.86 & 0.74 \\
fix-code-vulnerability & 10/10 & 10/10 & 10.78 & 4.26 & 1.08 & 0.43 \\
fix-git & 9/10 & 10/10 & 5.44 & 1.95 & 0.60 & 0.20 \\
fix-ocaml-gc & 9/10 & 10/10 & 23.21 & 28.09 & 2.58 & 2.81 \\
git-leak-recovery & 10/10 & 10/10 & 8.27 & 1.82 & 0.83 & 0.18 \\
log-summary-date-ranges & 10/10 & 10/10 & 5.03 & 3.24 & 0.50 & 0.32 \\
modernize-scientific-stack & 10/10 & 10/10 & 4.98 & 2.33 & 0.50 & 0.23 \\
openssl-selfsigned-cert & 10/10 & 10/10 & 5.50 & 2.35 & 0.55 & 0.24 \\
password-recovery & 10/10 & 10/10 & 19.66 & 10.46 & 1.97 & 1.05 \\
regex-log & 10/10 & 10/10 & 15.99 & 12.35 & 1.60 & 1.23 \\
\bottomrule
\end{tabular}
\end{table}

\clearpage
\bibliographystyle{plainnat}
\bibliography{../refs}

\end{document}
